\section{Discussão}

A análise e avaliação das estruturas de dados desenvolvidas permitiu chegar a conclusões claras face as perguntas propostas nos objetivos.
Relativamente à primeira pergunta, a implementação da Dense Advanced Myers List concluiu que não foi possível obter uma operação cons em
tempo constante. Esta limitação deve-se à impossibilidade de aceder as células, que são utilizadas como ponteiro uncle,
greater ou inc, sem a necessidade de atravessar a lista ou fazer um lookup. No entanto, foi possível implementar uma operação cons funcional 
para esta estrutura, com uma complexidade temporal de [INSERIR A COMPLEXIDADE]. Importa destacar que o objetivo de um cons com tempo constante 
exige uma restruturação da arquitetura, objetivo para futuras investigações.

Apesar de o cons em O(1) não ter sido obtido, a Dense Advanced Myers List apresentou resultados eficazes na leitura, tendo um tempo de execução
de lookup total de 60,35 ms, valor muito inferior à variante Advanced Myers List com 401,43 ms. Além disso, ainda apresentou um melhoramento no 
número de travessias, obtendo o melhor desempenho no pior caso quando comparada a todas as outras estruturas. Este ganho de performance tem um 
custo direto associado à memória exigindo 75,02 MB para um milhão de elementos, sendo a estrutura mais pesada do nosso estudo. Assim, conclui-se 
que esta estrutura é ideal para cenários onde a travessia de ponteiros é o principal obstáculo.

Relativamente à segunda pergunta, a implementação da IARASS revelou ser um sucesso na diminuição do problema da dupla descida. Ao introduzir células 
skip com um apontador adicional, conseguimos reduzir o número de travessias no pior caso de 35 (variante Improved) para 27. Além disso, apresentou um 
tempo de execução de 65,00 ms, ultrapassando a Improved Myers List (77,46 ms) e a Advanced Myers List (401,43 ms). No que toca a uso de memória, a 
IARASS utiliza 50,86 MB, valor superior aos 38,15 MB da variante Improved e consumindo menos memória do que a Advanced Myers List (53,41 MB). 
Assim, conclui-se que esta estrutura apresenta um melhor desempenho global superior ao da Advanced Myers List, otimizando simultaneamente a velocidade
de execução, o número de travessias e a eficiência de memória.

Em suma, a remoção da restrição permitiu-nos implementar duas novas estruturas de dados com propostas de valor distintas.
A Dense Advanced Myers List é a melhor escolha quando a velocidade lookup é o mais importante e a memória não é um problema a ser considerado, 
continuando com a operação de cons. Por sua vez, a IARASS mostra ser uma estrutura de dados equilibrada que oferece melhorias consideráveis no 
desempenho de lookup com um aumento moderado na memória.